\documentclass[11pt]{article}
\usepackage{fontspec}
\usepackage[french,english]{babel}
\usepackage{amsmath}
\usepackage{amssymb}
\usepackage{amsthm}
\usepackage{mathrsfs}
\usepackage{tikz-cd}
\usepackage{enumitem}

\newcommand{\printedpage}[1]{\marginpar{\footnotesize p.~#1}}
\newcommand{\msp}[1]{\marginpar{\footnotesize [#1]}}
\newcommand{\defeq}{\overset{\text{def}}{=}}
\newcommand{\fnm}[1]{\textsuperscript{(}\footnotemark[#1]\textsuperscript{)}}
\newcommand{\cl}{^{\tilde{}}}% tilde printed as exponent after a symbol or a parenthesis

\begin{document}

\printedpage{39}
% head: A. Grothendieck : Dérivateurs, Ch. XVII, Édition des Universités de Montpellier II et Paris VII
which decomposes into two
\[
\left\{
\begin{array}{rcl}
v_{!}(TC) & \subset & TC' \\
T\tilde{C}' & \subset & (v_{!}(TC))\cl.
\end{array}
\right.
\tag*{(4.7.3)}
\]
The less delicate of the two is the \msp{page 80} first, $v_{!}(TC) \subset TC'$, which amounts, by virtue of the definitions, implying
\[
v_{!}(TC) \subset v_{!}(C) \subset C' \, ,
\]
to
\[
v_{!}(TC) \;\subset\; W' \; ,
\]
i.e. by definition of $W'$, to
\[
v^{*}v_{!}(TC) \;\subset\; TC \; ,
\tag*{(4.7.4)}
\]
which is a hypothesis of type (4.6.1) $3^{\circ}$), to be added to the hypotheses (4.6.1).

The hypothetical relation
\[
\underbrace{T\tilde{C}'}_{=\,\tilde{C}'\cap W'} \;\subset\; (v_{!}(TC))\cl, \quad \text{or only} \quad TC' \subset (v_{!}(TC))\cl,
\tag*{$(*)$}
\]
appears to me to be of a distinctly more delicate nature. Since $T\tilde{C}' \supset TC' \supset v_{!}(TC)$, it is equivalent to the \textit{equality} of $(T\tilde{C}')_{*}$ resp. $(TC')_{*}$ and $(v_{!}(TC))_{*} = (v^{*})^{-1}(TC)_{*}$. Let us put
\[
\left\{
\begin{array}{lcl}
(TC)_{*} & \defeq & F \\
(TC')_{*} & \defeq & F' \;\supset\; (T\tilde{C}')_{*} \;=\; F'_{\sim} \; .
\end{array}
\right.
\tag*{(4.7.5)}
\]
We therefore have, by virtue of the inclusion
\[
v_{!}(TC) \;\subset\; TC' \; ( \;\subset\; T\tilde{C}' \; ) \; ,
\]
the relation
\[
F'_{\sim} \;\subset\; F' \;\subset\; (v^{*})^{-1}(F) \; ,
\tag*{(4.7.6)}
\]
and the question is therefore whether we have equality $F' = (v^{*})^{-1}(F)$, or even $F'_{\sim} = (v^{*})^{-1}(F)$ (whence $F'_{\sim} = F'$).

\msp{page 81}

In fact, we show fairly easily that we have $F'_{\sim} = F'$, i.e. $T\tilde{C}' \defeq \tilde{C}' \cap W' \subset (C' \cap W')\cl$ -- it is a general story about special categories with cofibrations \textsuperscript{(}\footnote[21]{cf. prop. 5.1 below.}\textsuperscript{)}. But I hardly see how to prove (under suitable hypotheses) that $F' = (v^{*})^{-1}(F)$. I shall have to come back to this if need be \textsuperscript{(}\footnote[22]{see however prop. 5.8 below.}\textsuperscript{)}. For the moment, I shall confine myself to reaping the fruit of 4.6 in a particular case:

\textbf{Proposition 4.8.} \textit{Let} $(\mathcal{M}, W, C)$ \textit{be a special category with cofibrations} (def. 4.4, page 94)\textit{. Let $I$ be an object of} Cat\textit{, $I_0$ the associated discrete category, let us consider $v : I_0 \longrightarrow I$ and} % [sic] page 94 -> page 74, apparently: the same déf. 4.4 is cited as "page 74" on printed p. 42
% foot: Version 9 février 2019
\printedpage{40}
% head: A. Grothendieck : Dérivateurs, Ch. XVII, Édition des Universités de Montpellier II et Paris VII
\[
\begin{tikzcd}
\mathcal{M}(I_0) \arrow[r, shift left, "v_{!}"] & \mathcal{M}(I) \arrow[l, shift left, "v^{*}"]
\end{tikzcd} \; ,
\]
\textit{finally let us endow $\mathcal{M}(I_0) = \mathcal{M}^{I_0}$ with the structure of a category with cofibrations,} product \textit{of those of the factors, i.e. let us define $W(I_0)$, $C(I_0)$ ``fiber by fiber''. (This coincides with the general definition given in paragraph} 2\textit{, for $I_0$ endowed with the discrete order structure, and makes $\mathcal{M}(I_0)$ a special category with cofibrations.) Then the couple $(v_{!}, v^{*})$ satisfies} \msp{page 82} \textit{the conditions of proposition} 4.6\textit{, as well as the condition supplementary to} (4.6.1)\textit{,}
\[
4^{\circ}) \qquad v^{*}(v_{!}(TC(I))) \;\subset\; TC(I) \; .
\tag*{(4.8.1)}
\]
\textit{Thus we find a structure of a special category with cofibrations on $\mathcal{M}(I)$, by putting}
\[
W(I) \;=\; (v^{*})^{-1}(W(I_0))
\tag*{(4.8.2)}
\]
\textit{(this is in fact equivalent to the usual definition of $W(I)$) and}
\[
C(I) \quad = \quad
\begin{array}[t]{l}
\textit{closure of } v_{!}(C(I_0)) \textit{ under the operations of cobase}\\
\textit{change and transfinite composition.}
\end{array}
\tag*{(4.8.3)}
\]
\textit{We have moreover under these conditions:}
\[
\underbrace{TF(I)}_{\defeq\,(C(I))_{*}} \;=\; (v^{*})^{-1}( \; \underbrace{TF(I_0)}_{\substack{\defeq\,(C_{I_0})_{*},\\ \textit{is also defined}\\ \textit{fiber by fiber}}} \; ) \; ,
\tag*{(4.8.8)}
\]
\[
\left\{
\begin{array}{rcl}
v_{!}( \, \overbrace{(TC)(I_0)}^{C(I_0)\cap W(I_0)} \, ) & \subset & TC(I) \qquad\qquad\qquad \text{\fnm{23}} \\
F(I) \defeq (TC(I))_{*} & \subset & (v^{*})^{-1}(F(I_0)).
\end{array}
\right.
\tag*{(4.8.9)}
\]
\footnotetext[23]{These inclusions should be equalities, but I do not know how to prove it without making stronger hypotheses, of the type \textit{special Quillen category.}}

\textsc{Proof}. We already know that $\mathcal{M}(I)$ is stable under small $\varinjlim$ , and that every element is accessible, as follows at once from the appropriate properties of $\mathcal{M}$. It is also clear that $v^{*}$ commutes with small $\varinjlim$ . It remains to prove (4.6.1. $3^{\circ}$), and $4^{\circ}$ in (4.8.1), i.e. \msp{page 83} that $v^{*}v_{!} : \mathcal{M}(I_0) \longrightarrow \mathcal{M}(I_0)$ maps $C$ into $C$, $TC$ into $TC$. Now we have
\[
v_{!}(F_0)(i) \;=\; \varinjlim_{(\alpha,j)\in i\backslash I_0} F_0(j) \qquad\quad \text{\fnm{24}}
\]
\footnotetext[24]{where $j \in \mathrm{Ob}\, I_0 = \mathrm{Ob}\, I$, and $\alpha : i \longrightarrow v(j) = j$.}
and $i\backslash I_0$ is the discrete category over $I_0$, whose fiber at $j \in \mathrm{Ob}\, I_0 \;=\; \mathrm{Ob}\, I$ is $\mathrm{Hom}_I(i,j)$:
\[
i\backslash I_0 \;\;=\;\; \coprod_{j\in \mathrm{Ob}\, I} \text{(discrete category associated to $\mathrm{Hom}_I(i,j)$)} \; ,
\]
% foot: Version 9 février 2019
\printedpage{41}
% head: A. Grothendieck : Dérivateurs, Ch. XVII, Édition des Universités de Montpellier II et Paris VII
thus
\[
v_{!}(F_0)(i) \;=\; \coprod_{j\in \mathrm{Ob}\, I = \mathrm{Ob}\, I_0} \underbrace{F_0(j)^{\mathrm{Hom}_I(i,j)}}_{\substack{\text{direct sum}\\ \text{of } \mathrm{Hom}_I(i,j)\\ \text{copies of } F_0(j)}} \, ,
\]
so we find
\[
v^{*}v_{!}F_0( \, \underbrace{\; i \;}_{\in \mathrm{Ob}\, I_0 \,=\, \mathrm{Ob}\, I} \, ) \;=\; v_{!}(F_0)(i) \;=\; \text{\tikz[baseline=0pt]\draw[->] (0,0) .. controls (0.4,0) .. (0.5,0.45);} \quad .
\]
% the curved arrow points up to the underbraced expression F_0(j)^{Hom_I(i,j)} of the preceding display
Since $C$ and $TC$ are stable under arbitrary direct sums (being so under finite direct sums, because of the finite stability conditions on $C$ and $TC$, and also under infinite sums because of the axiom of transfinite stability \dots\textsuperscript{(}\footnote[25]{deserves a general lemma \dots}\textsuperscript{)}), it does indeed follow that $v^{*}v_{!}$ maps $C(I_0)$ and $TC(I_0)$ into themselves. % [sic] somme infinies -> sommes infinies, agreement with the plural adjective

Thus, the first assertion of 4.8 follows from 4.6. Moreover, (4.8.8) and (4.8.9) were seen in 4.7, under the general conditions of 4.6 and assuming in addition the $4^{\circ}$ of (4.7.4).

\msp{page 84}

\section*{5\quad Relation of special categories with cofibrations to special Quillen categories}

Let $(\mathcal{M}, W, C)$ be a special category with cofibrations. I introduce
\[
TF \;=\; C_{*}, \qquad\quad TC \;=\; C \cap W, \qquad\quad F \;=\; (TC)_{*} \; ,
\tag*{(5.1)}
\]
so we have
\[
\begin{tikzcd}[row sep=large, column sep=large]
TC \arrow[r, hook] \arrow[d, leftrightarrow] & C \arrow[d, leftrightarrow, "\shortmid"{marking, sloped}] \\
F & TF \arrow[l, hook']
\end{tikzcd} \; ,
\tag*{(5.2)}
\]
$(C, TF)$ being a right closed Quillen precouple. We then know (XV 4.2 cor. 1, page 31) that $\tilde{C}$ consists of the direct factors of elements of $C$. Let us put
\[
T\tilde{C} \;=\; \tilde{C} \cap W \; ,
\tag*{(5.3)}
\]
whence
\[
TC \;\subset\; T\tilde{C} \; .
\tag*{(5.4)}
\]

\textbf{Proposition 5.1} \textit{Let $(\mathcal{M}, W, C)$ be a special category with cofibrations. We have $T\tilde{C} \subset (TC)\cl$ (and even $T\tilde{C} \subset (TC)^{\natural}$ \textsuperscript{(}\footnote[26]{\textbf{N.B.} For a subset $A$ of $\mathrm{Fl}(\mathcal{M})$, $A^{\natural}$ denotes the set of direct factors of elements of $A$.}\textsuperscript{)}), or what amounts to the same thing by virtue of (5.4), $TC$ and $T\tilde{C}$ have the same left $Q$-closure:}
\[
T\tilde{C} = (T\tilde{C})\cl = (TC)\cl \qquad\quad \textit{(and even } T\tilde{C} \subset (TC)^{\natural}\textit{),}
\tag*{(5.5)}
\]
% foot: Version 9 février 2019
\printedpage{42}
% head: A. Grothendieck : Dérivateurs, Ch. XVII, Édition des Universités de Montpellier II et Paris VII
\textit{whence}
\[
(\tilde{C})_{*} = C_{*} \, (= TF), \qquad\quad (T\tilde{C})_{*} = (TC)_{*} \, (= F) \qquad\quad \text{\fnm{27}}.
\tag*{(5.6)}
\]
\footnotetext[27]{These relations do not depend on the hypothesis $W$ stable under direct factors.}
\textit{Moreover, if $W$ stable under direct factors (or if merely every direct factor of an arrow of $TC$ is in $W$), and if moreover $W$ is stable under transfinite composition, then $(\mathcal{M}, W, \tilde{C})$ is a special category with cofibrations. Moreover, $T\tilde{C}$ is left $Q$-saturated, hence} (by 5.5)
\[
T\tilde{C} \;=\; (TC)^{\natural} \; .
\tag*{(5.7)}
\]

\msp{page 85}

\textsc{Proof}. Let us prove
\[
T\tilde{C} \;\subset\; (TC)\cl \, ,
\]
and more precisely, that every element of $T\tilde{C}$ is a direct factor of an element of $TC$:
\[
T\tilde{C} = (TC)^{\natural} \; .
\]
Let then $f \in T\tilde{C} \defeq \tilde{C}\cap W$, let us consider their factorization of $f$ as $pi$, with $i \in C$, $p \in TF$. Since $TF \subset W$, we have $p \in W$, hence $i \in W$, hence $i \in C \cap W \defeq TC$. Since $f \in C$, $p \in TF$, there exists $q : Y \longrightarrow X'$ such that $pq = \mathrm{id}_Y$, $qf = i$,
\[
\begin{tikzcd}[column sep=small, row sep=small]
X \arrow[rr, "f\in C\cap W"] \arrow[dr, "i\in C"'] & & Y \arrow[dl, dotted, shift right, "q"'] \\
& X' \arrow[ur, shift right, "p\in TF"'] &
\end{tikzcd}
\]
so $f$ appears as a direct factor of $i \in TC$, q.e.d.

Let us prove that $(\mathcal{M}, W, \tilde{C})$ is a special category with cofibrations (def. 4.4, page 74). Condition a) on $\mathcal{M}$ settled, condition b), the part on $W$, too. The stability of $\tilde{C}$ under cobase change and under composition OK, also (for b$'$) the stability under transfinite composition. Stability of $T\tilde{C} = \tilde{C} \cap W$ under cobase change: it remains to prove that if we have a cocartesian square
\[
\begin{tikzcd}[row sep=large, column sep=large]
X \arrow[r, "i"] \arrow[d] \arrow[dr, phantom, "\text{cocart.}"] & Y \arrow[d] \\
X' \arrow[r, "i'"] & Y'
\end{tikzcd} \; ,
\]
if $i \in \tilde{C}\cap W$, then $i' \in W$ (for we already know that $i' \in \tilde{C}$). But the preceding proof shows that it is a direct factor \textit{in} $X\backslash\mathcal{M}$ of an arrow in $TC$, so $i'$ is a direct factor of an arrow of $TC \subset W$. So if $W$ is stable under direct factors, $i' \in W$.
% foot: Version 9 février 2019

\printedpage{43}
% head: A. Grothendieck : Dérivateurs, Ch. XVII, Édition des Universités de Montpellier II et Paris VII
\msp{page 86}

This finishes proving b). For b$'$), it remains to see that $T\tilde{C}$ is stable under transfinite composition. I see no other way of seeing it, than by assuming that $W$ itself is stable under transfinite composition (which seems scarcely less demanding than the stability of $W$ under $\varinjlim$ \texttt{[filtering]}). This finishes proving b, b$'$). The condition c) on $C$ implies the same condition on $\tilde{C}$. Finally, since $\tilde{C}_{*} = C_{*}$, the condition d) is the same for $C$ and $\tilde{C}$. This finishes proving that $(\mathcal{M}, W, \tilde{C})$ is a special category with cofibrations. (\textbf{N.B.} we needed two additional hypotheses on $W$, stability under direct factors and under transfinite composition, to establish that $T\tilde{C}$ is stable under cobase change and under transfinite composition). The last assertion of 5.1 and the relation (5.7) which results from it, amounts to saying that $T\tilde{C} = \tilde{C} \cap W$ is stable under direct factors, which is obvious by the same property of $\tilde{C}$, and the hypothesis made on $W$ of stability under direct factors (or only that the direct factors of elements of $T\tilde{C}$ are in $W$, hence in $T\tilde{C}$). % [sic] Le dernière -> La dernière, gender agreement with "assertion"

\msp{page 87}

\textbf{Proposition 5.2.} \textit{Let $(\mathcal{M}, W, C)$ be a special category with cofibrations. Then with the notations of} (5.1)\textit{, we have}
\[
W \;=\; TF \circ TC,
\tag*{(5.8)}
\]
\[
TF \;=\; F \cap W.
\tag*{(5.9)}
\]

\textsc{Proof}. Let $f \in W$, let us write $f = pi$, with $i \in C$, $p \in TF$, hence $p \in W$, hence $i \in W$, hence $i \in C \cap W = TC$, whence (5.8).

Let us prove (5.9). We have $TF \subset F \cap W$, let us prove conversely that an $f \in F \cap W$ is in $TF$. Let us write $f = pi$, with $p \in TF$, $i \in TC$.
\[
\begin{tikzcd}[column sep=small, row sep=small]
X \arrow[rr, "f\in F\cap W"] \arrow[dr, "i\in TC"'] & & Y \\
& X' \arrow[ur, "p\in TF"'] &
\end{tikzcd}
\]
Since $f \in F$, $i \in TC$, and $TC \longleftrightarrow F$ by definition of $F$, $\exists\, r : X' \longrightarrow X$ such that $ri = \mathrm{id}_X$, $fr = p$. So $f$ is a direct factor of $p \in TF$, and since $TF = C_{*}$ is stable under direct factors, we conclude that $f \in TF$.

\textbf{Corollary 5.3.} \textit{Let $(\mathcal{M}, W, C)$ be a special category with cofibrations. For $(W, C, F)$, $(F \defeq (TC)_{*})$ to be a Quillen triple in $\mathcal{M}$, it is necessary and sufficient that the couple $(TC, F)$ satisfy the factorization axiom (i.e. that it be a Quillen precouple). For this Quillen structure to be closed, it is necessary and sufficient that $C$ be stable under direct factors.} % [sic] soit clos -> soit close, agreement with "cette structure"

\msp{page 88}

\textsc{Proof}. The necessity of the factorization axiom for $(TC, F)$, in order to have a Quillen triple, is tautological. For sufficiency, already having the two Quillen precouples $(C, TF)$, $(TC, F)$ with
% foot: Version 9 février 2019

\end{document}
